
%%%%%%%%%%%%%%%%%%%%%%%%%%%%%%%%%%%%%%%%
%           main body
%%%%%%%%%%%%%%%%%%%%%%%%%%%%%%%%%%%%%%%%


%-----------------------------------
%	TITLE PAGE
%-----------------------------------

\title[Convex Optimization]{\LARGE \bfseries 第四部分\quad 凸优化} % The short title appears at the bottom of every slide, the full title is only on the title page
\author[Convex Optimization Problems] % Your institution as it will appear on the bottom of every slide, may be shorthand to save space
{\Large \bfseries \S 3. 凸优化问题}
% \author{Singular Value Decomposition} % Your name
% \date{\today} % Date, can be changed to a custom date
\date{}

%------------------------------------------------

\begin{document}


%------------------------------------------------

\begin{frame}

\titlepage % Print the title page as the first slide

\end{frame}

%------------------------------------------------

\begin{frame}
\frametitle{内容提要} % Table of contents slide, comment this block out to remove it
\tableofcontents % Throughout your presentation, if you choose to use \section{} and \subsection{} commands, these will automatically be printed on this slide as an overview of your presentation
\end{frame}

%-----------------------------------
%	PRESENTATION SLIDES
%-----------------------------------

\section{标准形式的凸优化问题}

%------------------------------------------------

\begin{frame}

\begin{block}{标准形式的{\color{cyan}凸}优化问题}
\begin{eqnarray*}
& \text{minimize} & f_0(\mathbf{x}) \\
& \text{subject to} & f_i(\mathbf{x}) \le 0, \quad i = 1,\ldots,m \\
& & \mathbf{a}_j^T\mathbf{x} = b_j, \quad j = 1,\ldots,p
\end{eqnarray*}
\begin{itemize}
\item $\mathbf{x} \in \mathbb{R}^n$被称作优化变量
\item {\color{cyan}凸函数}$f_0$被称作目标函数
\item {\color{cyan}凸函数}$f_i$被称作不等式约束函数
\item {\color{cyan}仿射映射}$\mathbf{a}_j^T\mathbf{x} = b_j$被称作等式约束，等式约束也可以写作是$A\mathbf{x} = \mathbf{b}$，或是一般问题下的形式$h_j(\mathbf{x}) = 0$
\end{itemize}

\vspace{1em}

如果没有约束函数（即$m=p=0$），则称为无约束问题。
\end{block}

\end{frame}

%------------------------------------------------

\begin{frame}

\begin{block}{全局最优与局部最优}
全局最优解指的是以下的
$$\mathbf{x}^* = \argmin\limits_{\mathbf{x} \in \mathcal{D}} \left\{ f_0(\mathbf{x}) \ \middle|\ f_i(\mathbf{x}) \le 0, \; \mathbf{a}_j\mathbf{x} = b_j \right\}$$

局部最优解指的是如下的量
$$\mathbf{x}^* = \argmin\limits_{\mathbf{x} \in \mathcal{D}} \left\{ f_0(\mathbf{x}) \ \middle|\ f_i(\mathbf{x}) \le 0, \; \mathbf{a}_j\mathbf{x} = b_j, \; {\color{red} \Vert \mathbf{x} - \mathbf{x}^* \Vert \le R}\right\}$$

\vspace{1em}
\pause

\begin{itemize}
\item 对于凸优化问题来说，局部最优解就是全局最优解。
\end{itemize}
\end{block}

\end{frame}

%------------------------------------------------

\begin{frame}

\begin{block}{目标函数可微时的最优性准则}
记$X$为凸优化问题的可行集，即
$$X = \{ \mathbf{x} \in \mathcal{D} \ |\ f_i(\mathbf{x}) \le 0, \; i = 1,\ldots,m;\; \mathbf{a}_j\mathbf{x} = b_j, \; j=1,\ldots,p \}$$
那么$X$中的元素$\mathbf{x}$是最优解，当且仅当
$$\nabla f_0(\mathbf{x})^T(\widetilde{\mathbf{x}} - \mathbf{x}) \ge 0$$
对所有的$\widetilde{\mathbf{x}} \in X$成立
\end{block}

\end{frame}

%------------------------------------------------

\section{等价凸优化问题}

%------------------------------------------------

\begin{frame}

\begin{block}{等价凸优化问题}
两个凸优化问题被称作是等价的（Equivalent），如果从一个问题的最优解，很容易就能得到另一个问题的最优解，并且反之亦然。
\end{block}

\vspace{1em}

例如，将标准的凸优化问题进行伸缩
\begin{eqnarray*}
& \text{minimize} & \alpha_0f_0(\mathbf{x}) \\
& \text{subject to} & \alpha_i f_i(\mathbf{x}) \le 0, \quad i = 1,\ldots,m \\
& & \beta_j(\mathbf{a}_j^T\mathbf{x} - b_j) = 0, \quad j = 1,\ldots,p
\end{eqnarray*}
其中$\alpha_i > 0, \beta_j\neq 0$, 就与原来的标准的凸优化问题是等价的问题。
\end{frame}

%------------------------------------------------

\begin{frame}

\begin{block}{常见的等价凸优化问题}
\begin{itemize}
\item 消去等式约束：
\begin{eqnarray*}
& \text{minimize} & f_0(\mathbf{x}) \\
& \text{subject to} & f_i(\mathbf{x}) \le 0, \quad i = 1,\ldots,m \\
& & A\mathbf{x} = \mathbf{b}
\end{eqnarray*}
等价于如下的关于$\mathbf{z}$的凸优化问题
\begin{eqnarray*}
& \text{minimize} & f_0(F\mathbf{z} + \mathbf{x}_0) \\
& \text{subject to} & f_i(F\mathbf{z} + \mathbf{x}_0) \le 0, \quad i = 1,\ldots,m
\end{eqnarray*}
其中$\mathbf{x} = F\mathbf{z} + \mathbf{x}_0$是线性方程组$A\mathbf{x} = \mathbf{b}$的通解。
\end{itemize}
\end{block}

\end{frame}

%------------------------------------------------

\begin{frame}

\begin{block}{常见的等价凸优化问题}
\begin{itemize}
\item 添加等式约束：
\begin{eqnarray*}
& \text{minimize} & f_0(A_0\mathbf{x} + \mathbf{b}_0) \\
& \text{subject to} & f_i(A_i\mathbf{x} + \mathbf{b}_i) \le 0, \quad i = 1,\ldots,m
\end{eqnarray*}
等价于如下的关于$(\mathbf{x},\mathbf{y}_i)$的凸优化问题
\begin{eqnarray*}
& \text{minimize} & f_0(\mathbf{y}_0) \\
& \text{subject to} & f_i(\mathbf{y}_i) \le 0, \quad i = 1,\ldots,m \\
& & \mathbf{y}_i = A_i\mathbf{x} + \mathbf{b}_i, \quad i = 0,\ldots,m
\end{eqnarray*}
\end{itemize}
\end{block}

\end{frame}

%------------------------------------------------

\begin{frame}

\begin{block}{常见的等价凸优化问题}
\begin{itemize}
\item 引入松弛变量（Slack Variables）：
\begin{eqnarray*}
& \text{minimize} & f_0(\mathbf{x}) \\
& \text{subject to} & \mathbf{a}_i^T\mathbf{x} - b_i \le 0, \quad i = 1,\ldots,m
\end{eqnarray*}
等价于如下的关于$(\mathbf{x},\mathbf{s})$的凸优化问题
\begin{eqnarray*}
& \text{minimize} & f_0(\mathbf{x}) \\
& \text{subject to} & \mathbf{a}_i^T\mathbf{x} - s_i = b_i, \quad i = 1,\ldots,m \\
& & s_i \le 0, \quad i = 1,\ldots,m
\end{eqnarray*}
\end{itemize}
\end{block}

\end{frame}

%------------------------------------------------

\begin{frame}

\begin{block}{常见的等价凸优化问题}
\begin{itemize}
\item 上境图形式：
\begin{eqnarray*}
& \text{minimize} & f_0(\mathbf{x}) \\
& \text{subject to} & f_i(\mathbf{x}) \le 0, \quad i = 1,\ldots,m \\
& & \mathbf{a}_j^T\mathbf{x} = b_j, \quad j = 1,\ldots,p
\end{eqnarray*}
等价于如下的关于$(\mathbf{x},t)$的被称作上境图形式的凸优化问题
\begin{eqnarray*}
& \text{minimize} & t \\
& \text{subject to} & f_0(\mathbf{x}) - t \le 0, f_i(\mathbf{x}) \le 0, \quad i = 1,\ldots,m \\
& & \mathbf{a}_j^T\mathbf{x} = b_j, \quad j = 1,\ldots,p
\end{eqnarray*}
\end{itemize}
\end{block}

\end{frame}

%------------------------------------------------

\section{常见凸优化问题}

%------------------------------------------------

\begin{frame}

\begin{block}{线性规划，Linear Programming, LP}
\begin{eqnarray*}
& \text{minimize} & \mathbf{c}^T \mathbf{x} + d \\
& \text{subject to} & \mathbf{f}_i^T \mathbf{x} \le g_i, \; i = 1,\ldots,m \\
& & A\mathbf{x} = \mathbf{b}
\end{eqnarray*}

\begin{itemize}
\item 注：对于不等式约束$\mathbf{f}_i^T \mathbf{x} \le g_i, \; i = 1,\ldots,m$, 我们通常将其简记作
$$F\mathbf{x} \preccurlyeq \mathbf{g}$$
其中矩阵$F = (\mathbf{f}_1^T, \ldots, \mathbf{f}_m^T)^T, \mathbf{g} = (g_1, \ldots, g_m)^T$
\end{itemize}
\end{block}

\end{frame}

%------------------------------------------------

\begin{frame}

\begin{block}{二次规划，Quadratic Programming, QP}
\begin{eqnarray*}
& \text{minimize} & \dfrac{1}{2}\mathbf{x}^T P  \mathbf{x} + \mathbf{q}^T \mathbf{x} + r \\
& \text{subject to} & C \mathbf{x} \preccurlyeq \mathbf{d} \\
& & A\mathbf{x} = \mathbf{b}
\end{eqnarray*}
其中$P$为半正定实对称阵
\end{block}

\end{frame}

%------------------------------------------------

\begin{frame}

\begin{block}{二阶锥规划，Second-Order Cone Programming, SOCP}
\begin{eqnarray*}
& \text{minimize} & \mathbf{q}^T \mathbf{x} \\
& \text{subject to} & \Vert A_i\mathbf{x} + \mathbf{b}_i \Vert_2 \le \mathbf{c}_i^T \mathbf{x} + d_i, \; i = 1,\ldots,m \\
& & F\mathbf{x} = \mathbf{g}
\end{eqnarray*}
其中的不等式约束条件被称作是二阶锥约束
\end{block}

\end{frame}

%------------------------------------------------

% \begin{frame}

% \begin{block}{Vandermonde行列式的计算}
% \begin{enumerate}
% % \addtocounter{enumi}{0}
% \item 把第$n-1$行乘以$-a_1$加到第$n$行，再把第$n-2$行乘以$-a_1$加到第$n-1$行。按照这种顺序依次向上，直到用第$2$行乘以$-a_1$加到第$1$行。由行列式性质，$V_n$的值不变。
% \[
% V_n = \det \begin{pmatrix}
% 1 & 1 & \cdots & 1 \\
% 0 & a_2-a_1 & \cdots & a_n-a_1 \\
% \vdots & \vdots & & \vdots \\
% 0 & a_2^{n-2}(a_2-a_1) & \cdots & a_n^{n-2}(a_n-a_1) \\
% \end{pmatrix}
% \]
% \end{enumerate}
% \end{block}

% \end{frame}

%------------------------------------------------
% % closing page
% \begin{frame}
% \HUGE{\centerline{\bfseries {\color{gray} The} {\color{zkblue} End}}}

% \pause

% \vspace{1.2em}

% \Large{\centerline{\bfseries {\color{gray}Thank} \quad {\color{zkblue} You}}}

% \end{frame}

%----------------------------------------------------------------------------------------

\end{document}
